\conclusion

В настоящей выпускной квалификационной работе была поставлена цель
разработать библиотеку на языке~C++, реализующую пул потоков с механизмом
work-stealing, и экспериментально исследовать её производительность в
сравнении с классическим пулом потоков при различных сценариях нагрузки.
Поставленная цель достигнута, все задачи решены. Ниже приведены основные
результаты работы.

В ходе анализа предметной области установлено, что переход к многоядерным
архитектурам сделал параллелизм неотъемлемым требованием к современному
программному обеспечению. Классические подходы к управлению потоками ---
модель «поток на задачу» и пул потоков с единой глобальной очередью ---
обнаруживают существенные ограничения с точки зрения масштабируемости
и эффективности балансировки нагрузки при большом числе ядер. Механизм
work-stealing обладает строгим теоретическим обоснованием и доказанной
оптимальностью по времени выполнения и потреблению памяти, однако
существующие промышленные реализации~--- Intel~TBB и Microsoft~PPL~---
не допускают лёгкого изучения и модификации, а стандартная библиотека~C++
не предоставляет готовой реализации пула потоков с work-stealing. Данное
обстоятельство обусловило актуальность настоящего исследования.

В ходе теоретического обоснования и проектирования решения формальная
модель параллельных вычислений на основе ориентированного ациклического
графа позволила строго обосновать эффективность механизма work-stealing.
Доказанная Блюмофе и Лейзерсоном оценка $T_P \leq T_1 / P + O(T_\infty)$
показывает, что алгоритм work-stealing обеспечивает время выполнения,
близкое к теоретическому минимуму, без каких-либо априорных знаний
о структуре вычислений. Закон Амдала определил ключевое требование
к реализации: накладные расходы на управление очередями задач должны
быть минимизированы, поскольку они увеличивают эффективную долю
последовательной части вычисления и ограничивают достижимое ускорение.
Спроектированная архитектура библиотеки включает два модуля с идентичным
интерфейсом метода \texttt{submit()}: классический пул потоков
\texttt{SimpleThreadPool} и пул потоков с механизмом work-stealing
\texttt{WorkStealingPool}. Идентичность интерфейса обеспечивает
взаимозаменяемость реализаций и корректность сравнительного исследования.
Выбор упрощённой реализации \texttt{WorkStealingDeque} на основе
\texttt{std::deque} и \texttt{std::mutex} вместо неблокирующего алгоритма
Чейза-Лева является обоснованным компромиссом: корректная реализация
неблокирующей очереди требует применения механизмов безопасного
освобождения памяти в конкурентной среде, реализация которых выходила
за рамки настоящей работы.

В ходе программной реализации и экспериментального исследования
был получен ряд практически значимых результатов. Практическая реализация
библиотеки потребовала решения нетривиальных проблем, не очевидных
на этапе проектирования: гонка данных в деструкторе
\texttt{SimpleThreadPool}, переполнение беззнакового типа в методе
\texttt{pop()} и трудноустранимые ошибки в lock-free реализации
\texttt{WorkStealingDeque} были выявлены с помощью инструмента
Thread~Sanitizer и исправлены до проведения измерений производительности.
Все~28 автоматических тестов корректности успешно пройдены, гонок данных
не обнаружено.

Результаты экспериментального исследования не подтвердили исходную
гипотезу о превосходстве \texttt{WorkStealingPool} по показателю
реального времени выполнения. Основной причиной является использование
блокирующей синхронизации в \texttt{WorkStealingDeque}, что нивелирует
теоретическое преимущество локальных очередей и подчёркивает
принципиальную важность корректной реализации неблокирующей двусторонней
очереди для достижения теоретически предсказанной эффективности механизма
work-stealing. Вместе с тем исследование выявило фундаментальный компромисс
между латентностью и энергоэффективностью двух реализаций.
\texttt{SimpleThreadPool} обеспечивает меньшую латентность реакции
на новые задачи, что делает его предпочтительным для систем с жёсткими
временными требованиями. \texttt{WorkStealingPool} потребляет в~3,8~раза
меньше процессорного времени при неравномерной нагрузке~--- при
16~потоках 2,75~мс против 10,4~мс~--- что делает его предпочтительным
для фоновых вычислений в системах с ограниченным энергопотреблением.
Таким образом, выбор реализации пула потоков должен определяться
конкретными требованиями приложения, а не общими соображениями
о теоретической эффективности того или иного алгоритма.

Перспективами дальнейшего развития данной работы являются следующие
направления. Во-первых, реализация неблокирующей двусторонней очереди
по алгоритму Чейза-Лева с применением механизма hazard~pointers
или эпохального удаления позволит устранить главное ограничение текущей
реализации и проверить теоретические гарантии алгоритма work-stealing
в условиях отсутствия блокирующей синхронизации. Во-вторых,
исследование влияния топологии памяти NUMA на производительность
планировщика представляет практический интерес для многопроцессорных
серверных систем: привязка каждого рабочего потока к конкретному
NUMA-узлу и ограничение кражи задач в пределах одного узла способны
существенно снизить латентность доступа к памяти. В-третьих,
адаптивное управление числом рабочих потоков в зависимости от
текущей загрузки системы позволит библиотеке автоматически
подстраиваться под изменяющийся характер нагрузки без ручной
настройки параметров. Наконец, расширение набора сценариев нагрузки
за счёт рекурсивных алгоритмов типа «разделяй и властвуй»~--- сортировки
слиянием, параллельного обхода деревьев~--- позволит исследовать
механизм work-stealing в тех условиях, для которых он изначально
проектировался и в которых его теоретические преимущества должны
проявляться наиболее полно.